\chapter{Appendix}
\label{ch:appendix}

\section{Observation Builder}

Listing~\ref{lst:observation} shows the core branching logic of the
observation builder, from \texttt{autoscaler/\allowbreak \allowbreak src/\allowbreak \allowbreak rl\_\allowbreak env/\allowbreak \allowbreak observation.py},
lines 192-233. This is the function responsible for turning raw Prometheus
and Kubernetes telemetry into the eleven-feature vector described in
Table~\ref{tab:observation-features}, including the missing-data handling
rules of Section~\ref{sec:invalid-telemetry}.

\begin{lstlisting}[language=Python,caption={Core observation-building logic (\texttt{autoscaler/\allowbreak \allowbreak src/\allowbreak \allowbreak rl\_\allowbreak env/\allowbreak \allowbreak observation.py}, lines 192-233).},label={lst:observation}]
def build_observation(
    metrics: RawMetricSnapshot,
    replicas: ReplicaState,
    seconds_since_last_effective_action: float,
    history: HistoryState,
    previous_normalized: Optional[List[float]],
    config: ObservationConfig = ObservationConfig(),
) -> ObservationBuildResult:
    validity = determine_validity(metrics)
    ready = replicas.ready

    if previous_normalized is None:
        # Episode-start fallback: raw-range midpoint (0.5 after normalization
        # for symmetric-around-midpoint ranges) per sec.7.
        previous_normalized = [0.5, 0.0, 0.5, 0.5, 0.5, 0.0, 0.5, 0.0, 0.0, 0.5, 0.5]

    out = list(previous_normalized)
    raw: dict = {name: None for name in FEATURE_NAMES}

    if not validity.sample_valid:
        # Structurally invalid step: hold every level/trend feature except
        # the two controller-tracked ones (ready_replica_norm always comes
        # straight from the Deployment's own replica-count fields ...)
        out[9] = _norm(float(ready), config.min_replicas, config.max_replicas)
        out[10] = _norm(seconds_since_last_effective_action, 0.0, config.seconds_since_scale_cap_s)
        return ObservationBuildResult(out, raw, validity, ready)

    cpu_raw, mem_raw, inprog_raw = per_replica_ratios(metrics, ready)
    latency_ms = effective_latency_ms(metrics, validity.latency_observation_state)
    request_rate = metrics.request_rate_rps
    error_rate = 0.0 if metrics.error_rate_rps is None else (
        metrics.error_rate_rps / request_rate if request_rate else 0.0
    )
    # ... (remaining normalization and EWMA-trend assignment omitted here;
    # full function continues in the project repository.)
\end{lstlisting}

\section{V2 Reward Implementation}

Listing~\ref{lst:reward-v2} shows \texttt{compute\_\allowbreak reward\_\allowbreak v2}, from
\texttt{autoscaler/\allowbreak \allowbreak src/\allowbreak \allowbreak rl\_\allowbreak env/\allowbreak \allowbreak reward\_\allowbreak v2.py}, lines 199-262, implementing
Equation~\ref{eq:v2-reward} and Table~\ref{tab:reward-components}.

\begin{lstlisting}[language=Python,caption={V2 reward computation (\texttt{autoscaler/\allowbreak \allowbreak src/\allowbreak \allowbreak rl\_\allowbreak env/\allowbreak \allowbreak reward\_\allowbreak v2.py}, lines 199-262).},label={lst:reward-v2}]
def compute_reward_v2(
    *,
    decision_state_valid: bool,
    outcome_state_valid: bool,
    p95_latency_ms: Optional[float],
    error_rate: Optional[float],
    deployment_current_replicas: Optional[int],
    cpu_per_replica: Optional[float],
    memory_per_replica: Optional[float],
    requested_delta: int,
    effective_delta: int,
    override_reason: str,
    is_reversal: bool,
    config: RewardConfigV2 = RewardConfigV2(),
) -> RewardResultV2:
    if not decision_state_valid:
        components = {
            "reliability": config.w_reliability * _clip01(error_rate) if error_rate is not None else 0.0,
            "sla": 0.0, "pressure": 0.0, "waste": 0.0,
            "action": 0.0, "reversal": 0.0, "boundary": 0.0,
        }
        total = -sum(v for k, v in components.items())
        return RewardResultV2(total, components, REWARD_VALIDITY_PARTIAL)

    action_side = _action_side_components(requested_delta, effective_delta, override_reason, is_reversal, config)

    if not outcome_state_valid:
        components = {
            "reliability": config.w_reliability * _clip01(error_rate) if error_rate is not None else 0.0,
            "sla": 0.0, "pressure": 0.0, "waste": 0.0,
            **action_side,
        }
        total = -sum(v for k, v in components.items())
        return RewardResultV2(total, components, REWARD_VALIDITY_PARTIAL)

    reliability = config.w_reliability * _clip01(error_rate if error_rate is not None else 0.0)

    excess = max(0.0, (p95_latency_ms - config.sla_p95_target_ms) / config.sla_p95_target_ms)
    sla = config.w_sla * _clip01(excess ** 2)

    pressure_cpu = _pressure(cpu_per_replica, config.cpu_safe_high, config.cpu_hard_limit)
    pressure_mem = _pressure(memory_per_replica, config.mem_safe_high, config.mem_hard_limit)
    pressure = config.w_pressure * max(pressure_cpu, pressure_mem)

    waste_applies = (
        deployment_current_replicas is not None
        and deployment_current_replicas > config.min_replicas
        and cpu_per_replica is not None and cpu_per_replica < config.cpu_waste_threshold
        and memory_per_replica is not None and memory_per_replica < config.mem_waste_threshold
        and p95_latency_ms < config.sla_p95_target_ms
        and (error_rate is None or error_rate == 0.0)
    )
    waste = config.w_waste if waste_applies else 0.0

    components = {
        "reliability": reliability, "sla": sla, "pressure": pressure, "waste": waste,
        **action_side,
    }
    total = -sum(v for k, v in components.items())
    return RewardResultV2(total, components, REWARD_VALIDITY_FULL)
\end{lstlisting}

\section{Safety Gating Logic}

Listing~\ref{lst:gating} shows \texttt{compute\_\allowbreak effective\_\allowbreak action} in full,
from \texttt{autoscaler/\allowbreak \allowbreak src/\allowbreak \allowbreak rl\_\allowbreak env/\allowbreak \allowbreak gating.py}, lines 43-77, implementing
the replica-bound clamp, the pending-scale gate, and the asymmetric
scale-down cooldown described in Section~\ref{sec:safety-layer}.

\begin{lstlisting}[language=Python,caption={Deterministic safety gating (\texttt{autoscaler/\allowbreak \allowbreak src/\allowbreak \allowbreak rl\_\allowbreak env/\allowbreak \allowbreak gating.py}, lines 43-77).},label={lst:gating}]
def compute_effective_action(
    requested_delta: int,
    current_desired: int,
    current_ready: int,
    seconds_since_last_effective_action: Optional[float],
    config: GatingConfig = GatingConfig(),
) -> GatingResult:
    if requested_delta == 0:
        return GatingResult(0, OVERRIDE_NONE, bound_clamped=False, pending_scale=False)

    candidate = current_desired + requested_delta
    if candidate < config.min_replicas or candidate > config.max_replicas:
        return GatingResult(0, OVERRIDE_REPLICA_BOUND_CLAMP, bound_clamped=True, pending_scale=False)

    if current_desired != current_ready:
        return GatingResult(0, OVERRIDE_PENDING_SCALE_GATE, bound_clamped=False, pending_scale=True)

    if requested_delta > 0:
        return GatingResult(requested_delta, OVERRIDE_NONE, bound_clamped=False, pending_scale=False)

    # requested_delta < 0: SCALE_DOWN, subject to the 60s cooldown.
    if (
        seconds_since_last_effective_action is not None
        and seconds_since_last_effective_action < config.scale_down_cooldown_seconds
    ):
        return GatingResult(0, OVERRIDE_SCALE_DOWN_COOLDOWN, bound_clamped=False, pending_scale=False)

    return GatingResult(requested_delta, OVERRIDE_NONE, bound_clamped=False, pending_scale=False)
\end{lstlisting}
\texttt{GatingConfig} is parameterized with \texttt{min\_\allowbreak replicas=2},
\texttt{max\_\allowbreak replicas=8} (V2), and \texttt{scale\_\allowbreak down\_\allowbreak cooldown\_\allowbreak seconds=60.0}.

\section{Safety-Aware Tie-Break Selector}

Listing~\ref{lst:tiebreak} shows \texttt{select\_\allowbreak safety\_\allowbreak aware\_\allowbreak action}
from \texttt{autoscaler/\allowbreak \allowbreak src/\allowbreak \allowbreak rl\_\allowbreak env/\allowbreak \allowbreak action\_\allowbreak selection.py}, lines 171-241,
implementing rules R1-R6 of Section~\ref{sec:tie-break} and
Figure~\ref{fig:tie-break-flow}.

\begin{lstlisting}[language=Python,caption={Safety-aware tie-break selector (\texttt{autoscaler/\allowbreak \allowbreak src/\allowbreak \allowbreak rl\_\allowbreak env/\allowbreak \allowbreak action\_\allowbreak selection.py}, lines 171-241).},label={lst:tiebreak}]
def select_safety_aware_action(
    probabilities: Sequence[float],
    epsilon: float = DEFAULT_NEAR_TIE_EPSILON,
    stressed: Optional[bool] = None,
    stress_reasons: Tuple[str, ...] = (),
) -> SafetyAwareActionResult:
    if epsilon < 0:
        raise ValueError(f"epsilon must be >= 0, got {epsilon}")

    probs = np.asarray(probabilities, dtype=float)
    if probs.shape != (EXPECTED_NUM_ACTIONS,):
        raise ValueError(f"expected exactly {EXPECTED_NUM_ACTIONS} probabilities, got shape {probs.shape}")
    if not np.all(np.isfinite(probs)):
        raise ValueError(f"probabilities must all be finite, got {probs.tolist()}")
    if np.any(probs < 0):
        raise ValueError(f"probabilities must be non-negative, got {probs.tolist()}")

    max_p = float(probs.max())
    within_epsilon = set(np.flatnonzero(probs >= max_p - epsilon).tolist())
    tie_set_names = tuple(ACTION_NAMES[i] for i in sorted(within_epsilon))

    if len(within_epsilon) == 1:
        chosen = next(iter(within_epsilon))
        return SafetyAwareActionResult(
            action_index=chosen, action_name=ACTION_NAMES[chosen], is_tie=False,
            probabilities=tuple(probs.tolist()), epsilon=epsilon, tie_set=tie_set_names,
            tie_rule_applied=TIE_RULE_CLEAR_WINNER, stressed=None, stress_reasons=(),
        )

    if DOWN_ACTION_INDEX in within_epsilon:
        rule = TIE_RULE_DOWN_INVOLVES_UP if UP_ACTION_INDEX in within_epsilon else TIE_RULE_DOWN_HOLD
        return SafetyAwareActionResult(
            action_index=HOLD_ACTION_INDEX, action_name=ACTION_NAMES[HOLD_ACTION_INDEX], is_tie=True,
            probabilities=tuple(probs.tolist()), epsilon=epsilon, tie_set=tie_set_names,
            tie_rule_applied=rule, stressed=None, stress_reasons=(),
        )

    # Only {HOLD, UP} remains (the only other possible 2-way tie with 3 actions).
    if stressed is None:
        return SafetyAwareActionResult(
            action_index=HOLD_ACTION_INDEX, action_name=ACTION_NAMES[HOLD_ACTION_INDEX], is_tie=True,
            probabilities=tuple(probs.tolist()), epsilon=epsilon, tie_set=tie_set_names,
            tie_rule_applied=TIE_RULE_UNKNOWN_STRESS_FALLBACK, stressed=None, stress_reasons=(),
        )

    chosen = UP_ACTION_INDEX if stressed else HOLD_ACTION_INDEX
    return SafetyAwareActionResult(
        action_index=chosen, action_name=ACTION_NAMES[chosen], is_tie=True,
        probabilities=tuple(probs.tolist()), epsilon=epsilon, tie_set=tie_set_names,
        tie_rule_applied=TIE_RULE_HOLD_UP_STRESS_DEPENDENT, stressed=stressed,
        stress_reasons=stress_reasons if stressed else (),
    )
\end{lstlisting}

\section{Trainer Session Loop}

Listing~\ref{lst:trainer} shows the setup portion of
\texttt{run\_\allowbreak one\_\allowbreak final\_\allowbreak session}, from
\texttt{scripts/\allowbreak \allowbreak train\_\allowbreak ppo\_\allowbreak final\_\allowbreak v2.py}, lines 604-693 (abridged), which
runs one 256-step training session: it resolves the session's attempt
directory, loads the V2 environment configuration, constructs the live
Kubernetes backend and the Gymnasium environment, and wraps it with an
episode ledger recorder before training begins.

\begin{lstlisting}[language=Python,caption={Training session setup (\texttt{scripts/\allowbreak \allowbreak train\_\allowbreak ppo\_\allowbreak final\_\allowbreak v2.py}, lines 604-693, abridged).},label={lst:trainer}]
def run_one_final_session(
    session_index: int,
    campaign_state: Dict[str, Any],
    seed: int = FINAL_TRAINING_SEED,
) -> Dict[str, Any]:
    from stable_baselines3 import PPO
    from autoscaler.src.rl_env.live_backend import LiveKubernetesBackend
    from collectors.metrics_collector import PrometheusClient
    from scripts.v2_pilot_safety import PROMETHEUS_URL, PilotSafetyCallback, get_pod_health
    from scripts.v2_traffic_orchestrator import V2TrafficOrchestrator, V2TrafficSyncedWrapper

    session_dir = resolve_session_attempt_dir(session_index)
    physical_attempt_number = physical_attempt_number_for_dir(session_dir)
    model_path = session_dir / SB3_CHECKPOINT_NAME
    rng_path = session_dir / RNG_STATE_NAME
    metadata_path = session_dir / SESSION_METADATA_NAME
    metrics_path = session_dir / SESSION_TRAINING_METRICS_NAME
    telemetry_path = session_dir / SESSION_TELEMETRY_NAME

    assert_final_namespace_isolated([model_path, FINAL_LEDGER_PATH])
    session_dir.mkdir(parents=True, exist_ok=False)

    env_config_v2 = load_env_config_v2_from_yaml()
    backend = LiveKubernetesBackend(
        prometheus_url=PROMETHEUS_URL,
        min_replicas=env_config_v2.gating.min_replicas,
        initial_replicas=env_config_v2.initial_replicas,
    )
    baseline_pods = get_pod_health("autoscaler-app", "app=demo-app")
    baseline_restart_counts = {p.name: p.restart_count for p in baseline_pods}

    raw_env = KubernetesAutoscalerEnvV2(
        backend=backend, episode_duration_s=600.0, config=env_config_v2,
    )
    controller_code_hash = hashlib.sha256(
        (PROJECT_ROOT / "autoscaler/src/rl_env/environment_v2.py").read_bytes()
        + (PROJECT_ROOT / "autoscaler/src/rl_env/reward_v2.py").read_bytes()
    ).hexdigest()
    config_hash = hashlib.sha256(
        (PROJECT_ROOT / "configs/rl_v2.yaml").read_bytes() + (PROJECT_ROOT / "configs/reward_v2.yaml").read_bytes()
    ).hexdigest()
    ledgered_env = LedgeredEpisodeWrapper(
        raw_env, ledger_path=FINAL_LEDGER_PATH, traffic_profile="PENDING_FIRST_RESET",
        profile_hash="PENDING_FIRST_RESET", seed=seed, controller_code_hash=controller_code_hash,
        config_hash=config_hash, role="v2_final_training",
    )
    # ... (environment reset, PPO.learn() call, checkpoint/metadata writing,
    # and the session's own post-hoc self-verification continue below; see
    # scripts/train_ppo_final_v2.py in the project repository.)
\end{lstlisting}

\section{Evaluation Runner Core}

Listing~\ref{lst:evalrunner} shows \texttt{run\_\allowbreak one\_\allowbreak attempt} in full, from
\texttt{scripts/\allowbreak \allowbreak evaluate\_\allowbreak v2\_\allowbreak campaign.py}, lines 392-460, the function
that runs exactly one evaluation or stress-test attempt end to end:
preflight, controller acquisition, traffic start/stop, the arm-specific
episode function, cleanup, and result persistence.

\begin{lstlisting}[language=Python,caption={Evaluation/stress attempt runner (\texttt{scripts/\allowbreak \allowbreak evaluate\_\allowbreak v2\_\allowbreak campaign.py}, lines 392-460).},label={lst:evalrunner}]
def run_one_attempt(run_record: RunRecord, ctx: CampaignContext, attempt_number: int) -> Dict[str, Any]:
    episode_id = f"{run_record.run_id}__attempt{attempt_number}"
    started_utc = datetime.now(timezone.utc).isoformat()
    entry: Dict[str, Any] = {
        "attempt_id": str(uuid.uuid4()), "run_id": run_record.run_id, "run_index": run_record.run_index,
        "profile": run_record.profile, "repeat_index": run_record.repeat_index, "arm_id": run_record.arm_id,
        "attempt_number": attempt_number, "episode_id": episode_id, "started_utc": started_utc,
    }
    handle = None
    try:
        preflight = ctx.preflight_fn()
        if not preflight.get("ok"):
            raise ToolingFailureError(f"preflight failed: {preflight}")

        if run_record.arm_id == "hpa_v2_cpu_only":
            resolve_cpu_only_controller_identity()

        profile_meta = ctx.profile_resolver_fn(run_record.profile)
        handle = ctx.orchestrator.start_episode_traffic(run_record.profile, episode_id)
        try:
            if run_record.arm_id == "ppo_v2":
                result = ctx.ppo_episode_fn(profile_meta, episode_id)
            elif run_record.arm_id == "hpa_v2_cpu_only":
                result = ctx.cpu_only_episode_fn(profile_meta, episode_id)
            else:
                result = ctx.hpa_episode_fn(profile_meta, episode_id)
        finally:
            ctx.orchestrator.stop_episode_traffic(handle, episode_id)
            ctx.orchestrator.drain_server_side_requests()

        cleanup_result = ctx.cleanup_fn()
        ctx.settle_fn(ctx.settle_seconds)
        entry.update({"status": "VALID", "result": result, "cleanup": cleanup_result})

    except CpuOnlyArmNotYetSupportedError as exc:
        cleanup_result = _safe_cleanup(ctx)
        entry.update({"status": "FAILED_CLOSED", "error": str(exc), "cleanup": cleanup_result})
    except SafetyTerminationError as exc:
        cleanup_result = _safe_cleanup(ctx)
        entry.update({"status": "SAFETY_TERMINATION", "error": str(exc), "cleanup": cleanup_result})
    except ToolingFailureError as exc:
        cleanup_result = _safe_cleanup(ctx)
        entry.update({"status": "TOOLING_FAILURE", "error": str(exc), "cleanup": cleanup_result})
    except Exception as exc:  # noqa: BLE001
        cleanup_result = _safe_cleanup(ctx)
        entry.update({"status": "TOOLING_FAILURE", "error": f"unclassified: {exc}", "cleanup": cleanup_result})

    entry["ended_utc"] = datetime.now(timezone.utc).isoformat()
    run_dir = ctx.runs_dir / f"{run_record.run_id}__attempt{attempt_number}"
    run_dir.mkdir(parents=True, exist_ok=True)
    with open(run_dir / "attempt.json", "w", encoding="utf-8") as f:
        json.dump(entry, f, indent=2, default=str)

    if entry["status"] == "VALID" and isinstance(entry.get("result"), dict) and "samples" in entry["result"]:
        series_dict = episode_result_to_run_series_dict(entry["result"], run_record.run_id, run_record.repeat_index)
        with open(run_dir / "normalized_series.json", "w", encoding="utf-8") as f:
            json.dump(series_dict, f, indent=2, default=str)

    return entry
\end{lstlisting}

\section{Configuration File Excerpts}

Listing~\ref{lst:hpa-cm} and Listing~\ref{lst:hpa-cpu} show the two V2 HPA
manifests, \texttt{k8s/\allowbreak \allowbreak hpa/\allowbreak \allowbreak hpa\_\allowbreak v2\_\allowbreak cpu\_\allowbreak memory.yaml} and
\texttt{k8s/\allowbreak \allowbreak hpa/\allowbreak \allowbreak hpa\_\allowbreak v2\_\allowbreak cpu\_\allowbreak only\_\allowbreak min2.yaml}. Both share the same
replica bounds (2-8) and scaling behaviour (immediate scale-up, a 120-second
scale-down stabilization window), differing only in which metric(s) drive
the scaling decision -- exactly the isolation CORR-60/CORR-62 were designed
to produce.

\begin{lstlisting}[caption={\texttt{k8s/\allowbreak \allowbreak hpa/\allowbreak \allowbreak hpa\_\allowbreak v2\_\allowbreak cpu\_\allowbreak memory.yaml} -- the HPA CPU+Memory (resource baseline) arm.},label={lst:hpa-cm}]
apiVersion: autoscaling/v2
kind: HorizontalPodAutoscaler
metadata:
  name: demo-app-hpa-v2
spec:
  scaleTargetRef:
    apiVersion: apps/v1
    kind: Deployment
    name: demo-app
  minReplicas: 2
  maxReplicas: 8
  metrics:
    - type: Resource
      resource:
        name: cpu
        target:
          type: Utilization
          averageUtilization: 260
    - type: Resource
      resource:
        name: memory
        target:
          type: Utilization
          averageUtilization: 120
  behavior:
    scaleUp:
      stabilizationWindowSeconds: 0
      policies:
        - type: Pods
          value: 2
          periodSeconds: 15
    scaleDown:
      stabilizationWindowSeconds: 120
      policies:
        - type: Pods
          value: 1
          periodSeconds: 60
\end{lstlisting}

\begin{lstlisting}[caption={\texttt{k8s/\allowbreak \allowbreak hpa/\allowbreak \allowbreak hpa\_\allowbreak v2\_\allowbreak cpu\_\allowbreak only\_\allowbreak min2.yaml} -- the HPA CPU-only (CPU baseline) arm (CORR-62).},label={lst:hpa-cpu}]
apiVersion: autoscaling/v2
kind: HorizontalPodAutoscaler
metadata:
  name: demo-app-hpa-v2-cpu
spec:
  scaleTargetRef:
    apiVersion: apps/v1
    kind: Deployment
    name: demo-app
  minReplicas: 2
  maxReplicas: 8
  metrics:
    - type: Resource
      resource:
        name: cpu
        target:
          type: Utilization
          averageUtilization: 260
  behavior:
    scaleUp:
      stabilizationWindowSeconds: 0
      policies:
        - type: Pods
          value: 2
          periodSeconds: 15
    scaleDown:
      stabilizationWindowSeconds: 120
      policies:
        - type: Pods
          value: 1
          periodSeconds: 60
\end{lstlisting}

Listing~\ref{lst:reward-config} shows the key fields of
\texttt{configs/\allowbreak \allowbreak reward\_\allowbreak v2.yaml}, the frozen configuration underlying the
V2 reward of Equation~\ref{eq:v2-reward} and Table~\ref{tab:reward-components}.

\begin{lstlisting}[caption={Key fields of \texttt{configs/\allowbreak \allowbreak reward\_\allowbreak v2.yaml}.},label={lst:reward-config}]
reward_contract_version: 2
weights:
  w_reliability: 6.0
  w_sla: 4.0
  w_pressure: 2.0
  w_waste: 0.3
  w_action: 0.1
  w_reversal: 0.5
  w_boundary: 0.2
sla:
  p95_target_ms: 1000      # raised from 300ms, CORR-57
pressure:
  cpu_safe_high: 3.0
  cpu_hard_limit: 5.0
  memory_safe_high: 1.2
  memory_hard_limit: 4.0
waste:
  cpu_waste_threshold: 1.0   # raised from 0.3, CORR-57
  memory_waste_threshold: 1.0  # raised from 0.7, CORR-57
reversal_window_seconds: 120
\end{lstlisting}

Listing~\ref{lst:rl-config} shows the key fields of
\texttt{configs/\allowbreak \allowbreak rl\_\allowbreak v2.yaml}, including the V2 replica floor (CORR-55) and
the provisional tie-break epsilon later superseded, for evaluation
purposes only, by the calibrated value of Section~\ref{sec:tie-break}.

\begin{lstlisting}[caption={Key fields of \texttt{configs/\allowbreak \allowbreak rl\_\allowbreak v2.yaml}.},label={lst:rl-config}]
rl_contract_version: 2
action_map:
  0: -1   # SCALE_DOWN_ONE
  1: 0    # HOLD
  2: 1    # SCALE_UP_ONE
ppo:
  learning_rate: 0.0003
  n_steps: 2048
  batch_size: 64
  n_epochs: 10
  gamma: 0.99
  gae_lambda: 0.95
  clip_range: 0.2
  ent_coef: 0.01
  vf_coef: 0.5
  net_arch: [128, 128]
scaling_v2:
  min_replicas: 2     # CORR-55 (V2 floor raised from 1)
  max_replicas: 8
  initial_replicas: 2
safe_action:
  near_tie_epsilon: 0.05
  epsilon_status: provisional_unvalidated
\end{lstlisting}

\section{Correction Log Summary}

Table~\ref{tab:corr-summary} lists every entry of the project's append-only
design-correction log (\texttt{docs/\allowbreak \allowbreak experimental\_\allowbreak design\_\allowbreak corrections.md}),
CORR-1 through CORR-65, one line per entry. Entries CORR-49 through CORR-54
are a known, acknowledged documentation gap: the correction log's own
CORR-55 entry states that these numbers were consumed in code (for
example, \texttt{scripts/\allowbreak \allowbreak run\_\allowbreak dynamics\_\allowbreak calibration\_\allowbreak slot.py}) but were
never written up as entries in this document; we report this gap rather
than inventing content for it. CORR-1 through CORR-26 concern the
project's pre-V2 calibration and HPA-tuning phases (Minikube resource
accounting, workload calibration, and the counterbalanced study that
selected the 260\% CPU target used by both V2 HPA baselines); CORR-27
onward concern the frozen MDP contract and the V2 multi-metric controller
redesign discussed throughout Chapter~\ref{ch:methodology} and
Chapter~\ref{ch:results}.

\begingroup
\small\singlespacing
\begin{longtable}{p{1.6cm}p{11.5cm}}
\caption{Correction log summary (CORR-1 through CORR-65).}
\label{tab:corr-summary}\\
\toprule
Entry & Title \\
\midrule
\endfirsthead
\toprule
Entry & Title \\
\midrule
\endhead
CORR-1 & Minikube Resource-Accounting Mismatch \\
CORR-2 & HPA Metric-Resolution Fairness for Benchmark \\
CORR-3 & CPU Limits and maxReplicas (Phase 6A correction) \\
CORR-4 & RL Agent Metric Source \\
CORR-5 & Workload Generation Isolation (Phase 6C Diagnostic Run Invalidation) \\
CORR-6 & Application Metric Label Collision with Prometheus Target Labels \\
CORR-7 & Status of \texttt{demo\_\allowbreak app\_\allowbreak last\_\allowbreak work\_\allowbreak intensity} Diagnostic Gauge \\
CORR-8 & Host Resource Isolation, Pod Lifecycle Resilience, and Cluster Headroom Guidelines \\
CORR-9 & Scientific Telemetry Collection Integrity, Target Health, and Runtime Provenance \\
CORR-10 & Workload Computation Reality, Host Sharing, and Planned Schedule Replay \\
CORR-11 & Dispatch Lag vs.\ Latency Attribution and Profile Duration Boundaries \\
CORR-12 & Stage-1 Calibration Artifact Validity, CPU Rate Terminology, and Throttling Definitions \\
CORR-13 & Canonical Workload Window and Telemetry Summary Methodology \\
CORR-14 & Empirical Capacity Boundaries, Workload-Window Dispatch Rate, and CFS Throttling Deltas \\
CORR-15 & Single-Replica Capacity Calibration Closure and Capacity Statement \\
CORR-16 & Initial HPA Baseline Implementation and Metrics Server Research Configuration \\
CORR-17 & HPA Tuning Methodology, Exact-Replay Capability, and Precommitted Selection Criteria \\
CORR-18 & HPA Tuning Methodology Hardening, Repetition/Order Protocol, Numeric SLA \\
CORR-19 & Event Baselining, Evidence Hierarchy, and Conditional 6-RPS Interpretation \\
CORR-20 & Corrected HPA Average-Utilization Semantics; Reconciliation of First Official 220\% Trial \\
CORR-21 & Corrected Multi-Replica Traffic-Delivery Path (\texttt{kubectl port-forward} Does Not Load-Balance) \\
CORR-22 & Eliminated Duplicate Prometheus Scraping; Fresh-Connection Discriminator Fix \\
CORR-23 & Fine-Grained Final-Path HPA-Relevant CPU Map \\
CORR-24 & Redesigned HPA Candidate Family (200/230/260) and Canonical Tuning Schedule v2 \\
CORR-25 & Standardized Exact-Interval Replica-Trajectory Integration \\
CORR-26 & Official HPA Tuning Dataset Closed; 260\% Selected as the Frozen Baseline \\
CORR-27 & RL MDP Contract V2: CPU-Clip, HPA-Coupling, Gating, and Held-Out-Leakage Corrections \\
CORR-28 & RL Runtime Dependency Compatibility and Pre-Environment Readiness Freeze \\
CORR-29 & RL Gymnasium Environment Mechanics and Pluggable-Backend Contract \\
CORR-30 & RL Transition-Reward Validity Decomposition and Deterministic Episode-Reset Control Age \\
CORR-31 & Empirical RL Dynamics Calibration Protocol Freeze \\
CORR-32 & RL Empirical-Dynamics Calibration Protocol Integrity Corrections \\
CORR-33 & Final RL Dynamics-Calibration Execution-Readiness Clarifications \\
CORR-34 & Global Calibration Run Order and Exhaustive Routing-Integrity Classification \\
CORR-35 & Exposure-Weighted Routing Opportunity and Transition-Safe Per-Pod Integrity Analysis \\
CORR-36 & Steady-State Ready-Count Invariant Can Be Violated by Overload-Induced Readiness Loss \\
CORR-37 & Post-First-Run Separation of Execution Validity and Fitting Eligibility \\
CORR-38 & Post-Data Expansion to Complete Relevant-Pod Observability \\
CORR-39 & Unverified External Workload-Delivery-Path Integrity During a Sustained Connection-Error Block \\
CORR-40 & Stable Host-to-Kubernetes Workload Transport Transition \\
CORR-41 & Kubernetes-Dataplane Workload-Ingress Evidence \\
CORR-42 & Temporal Dataplane-Evidence Acquisition Coverage \\
CORR-43 & Deterministic Short-Trailing-Window Evidence Partitioning \\
CORR-44 & Replica-Transition Materiality/Scored-Window Operationalization \\
CORR-45 & Replica-Transition Scale-Down Lag Target-State Operationalization \\
CORR-46 & Transition Collector Transient-Timeout Resilience \\
CORR-47 & Workload-Transition Pre-First-Use Operationalization \\
CORR-48 & Workload-Transition Trajectory Traffic-Window Boundary \\
CORR-49--54 & \textit{Documentation gap -- consumed in code but never written up as log entries (acknowledged in CORR-55).} \\
CORR-55 & V2 Replica Floor Raised 1 $\to$ 2 for Both V2 Controllers \\
CORR-56 & Idle (Zero-Traffic) Samples Incorrectly Classified Invalid \\
CORR-57 & V2 SLA Target and Waste Thresholds Recalibrated \\
CORR-58 & Safety-Aware Asymmetric Tie-Breaking for V2 Deterministic Inference \\
CORR-59 & Pre-Registered Budget Amendment, 5{,}120 $\to$ 12{,}288 Timesteps \\
CORR-60 & V2 Evaluation Pre-Registration \\
CORR-61 & Evaluation Safety/Termination Clarification \\
CORR-62 & CPU-Only HPA Arm Added (Post-Training) \\
CORR-63 & CPU-Only HPA Arm Blocked by Leftover \texttt{demo-app-hpa-v2} (Code Defect) \\
CORR-64 & \texttt{restore\_\allowbreak canonical\_\allowbreak idle\_\allowbreak state()} Left \texttt{demo-app-hpa-v2-cpu} Behind (Code Defect) \\
CORR-65 & Exploratory Post-Hoc Stress Test at Training-Level Peaks \\
\bottomrule
\end{longtable}
\endgroup

\section{Supplementary Data Tables}

Table~\ref{tab:training-summary-appendix} gives the full training-summary
statistics underlying the discussion of the V2 training campaign in
Chapter~\ref{ch:results}, including minimum and maximum values across all
48 sessions, not only the first/last values discussed in the main text.

\begin{table}[H]
\centering
\caption{V2 training summary statistics (48 analyzed sessions).}
\label{tab:training-summary-appendix}
\begin{tabular}{ll}
\hline
field & value \\
\hline
n\_sessions\_analyzed & 48 \\
first\_session\_entropy\_loss & -1.0984 \\
last\_session\_entropy\_loss & -0.2792 \\
min\_entropy\_loss & -1.0984 \\
max\_entropy\_loss & -0.2048 \\
first\_session\_explained\_variance & 0.0125 \\
last\_session\_explained\_variance & 0.7993 \\
min\_explained\_variance & -1.2461 \\
max\_explained\_variance & 0.7993 \\
n\_failed\_retried\_attempts & 2 \\
failed\_attempt\_session\_indices & 4, 38 \\
\hline
\end{tabular}

\end{table}

Table~\ref{tab:eval-secondary-appendix} gives the full secondary-metrics
table for the pre-registered \texttt{eval\_\allowbreak final} campaign, and
Table~\ref{tab:stress-descriptive-appendix} the full descriptive table for
the exploratory \texttt{stress\_\allowbreak explore} campaign, both reproduced here in
landscape orientation because of their column count.

\begin{landscape}
\begin{table}[H]
\centering
\caption{\texttt{eval\_\allowbreak final} secondary metrics, mean $\pm$ standard deviation over 3 repeats per (profile, arm) cell.}
\label{tab:eval-secondary-appendix}
\tiny
\resizebox{\linewidth}{!}{\begin{tabular}{lllllllllllll}
\hline
profile & arm & n\_repeats & mean\_p95\_latency\_ms & median\_p95\_latency\_ms & error\_rate & max\_per\_pod\_memory\_mib & scaling\_actions\_count & reversals\_count & reaction\_time\_s & pod\_restarts\_count & invalid\_samples\_count & sla\_violation\_rate\_300ms\_reference \\
\hline
CPU step & HPA CPU+Memory (resource baseline) & 3 & 270.9436 +/- 36.3194 & 232.6063 +/- 2.8184 & 0.0000 +/- 0.0000 & 48.5469 +/- 9.6242 & 0.0000 +/- 0.0000 & 0.0000 +/- 0.0000 & n/a & 0.0000 +/- 0.0000 & 0.0000 +/- 0.0000 & 0.1083 +/- 0.0520 \\
CPU step & HPA CPU-only (CPU baseline) & 3 & 249.6525 +/- 4.1099 & 231.9082 +/- 2.3202 & 0.0000 +/- 0.0000 & 56.7891 +/- 3.9332 & 0.0000 +/- 0.0000 & 0.0000 +/- 0.0000 & n/a & 0.0000 +/- 0.0000 & 0.3333 +/- 0.5774 & 0.0671 +/- 0.0137 \\
CPU step & PPO (multi-metric, 11 signals) & 3 & 274.1151 +/- 92.5713 & 229.3114 +/- 4.9961 & 0.0000 +/- 0.0000 & 51.6029 +/- 10.2576 & 2.0000 +/- 0.0000 & 0.0000 +/- 0.0000 & n/a & 0.0000 +/- 0.0000 & 0.0000 +/- 0.0000 & 0.0750 +/- 0.1299 \\
memory ramp & HPA CPU+Memory (resource baseline) & 3 & 975.0000 +/- 0.0000 & 975.0000 +/- 0.0000 & 0.0000 +/- 0.0000 & 123.7539 +/- 10.8288 & 0.0000 +/- 0.0000 & 0.0000 +/- 0.0000 & n/a & 0.0000 +/- 0.0000 & 0.0000 +/- 0.0000 & 1.0000 +/- 0.0000 \\
memory ramp & HPA CPU-only (CPU baseline) & 3 & 975.0000 +/- 0.0000 & 975.0000 +/- 0.0000 & 0.0000 +/- 0.0000 & 140.8529 +/- 8.9892 & 0.0000 +/- 0.0000 & 0.0000 +/- 0.0000 & n/a & 0.0000 +/- 0.0000 & 0.0000 +/- 0.0000 & 1.0000 +/- 0.0000 \\
memory ramp & PPO (multi-metric, 11 signals) & 3 & 950.6250 +/- 0.0000 & 975.0000 +/- 0.0000 & 0.0000 +/- 0.0000 & 108.5924 +/- 9.4735 & 1.0000 +/- 0.0000 & 0.0000 +/- 0.0000 & n/a & 0.0000 +/- 0.0000 & 0.0000 +/- 0.0000 & 0.9750 +/- 0.0000 \\
mixed bursts & HPA CPU+Memory (resource baseline) & 3 & 975.0000 +/- 0.0000 & 975.0000 +/- 0.0000 & 0.0000 +/- 0.0000 & 97.5573 +/- 18.3602 & 0.0000 +/- 0.0000 & 0.0000 +/- 0.0000 & n/a & 0.0000 +/- 0.0000 & 0.3333 +/- 0.5774 & 1.0000 +/- 0.0000 \\
mixed bursts & HPA CPU-only (CPU baseline) & 3 & 975.0931 +/- 0.0851 & 975.0000 +/- 0.0000 & 0.0000 +/- 0.0000 & 118.9987 +/- 0.1583 & 0.0000 +/- 0.0000 & 0.0000 +/- 0.0000 & n/a & 0.0000 +/- 0.0000 & 0.0000 +/- 0.0000 & 1.0000 +/- 0.0000 \\
mixed bursts & PPO (multi-metric, 11 signals) & 3 & 950.6250 +/- 0.0000 & 975.0000 +/- 0.0000 & 0.0000 +/- 0.0000 & 108.2292 +/- 8.9803 & 1.3333 +/- 0.5774 & 0.0000 +/- 0.0000 & n/a & 0.0000 +/- 0.0000 & 0.0000 +/- 0.0000 & 0.9750 +/- 0.0000 \\
\hline
\end{tabular}
}
\end{table}
\end{landscape}

\begin{landscape}
\begin{table}[H]
\centering
\caption{\texttt{stress\_\allowbreak explore} descriptive metrics, 1 repeat per (profile, arm) cell -- exploratory, CORR-65, never subject to the pre-registered decision rule.}
\label{tab:stress-descriptive-appendix}
\tiny
\resizebox{\linewidth}{!}{\begin{tabular}{llllllllllllllll}
\hline
profile & arm & n\_repeats & exploratory & sla\_violation\_rate\_1000ms & replica\_minutes & mean\_p95\_latency\_ms & median\_p95\_latency\_ms & error\_rate & max\_per\_pod\_memory\_mib & scaling\_actions\_count & reversals\_count & reaction\_time\_s & pod\_restarts\_count & invalid\_samples\_count & sla\_violation\_rate\_300ms\_reference \\
\hline
CPU bursts & HPA CPU+Memory (resource baseline) & 1 & yes & 0.5263 & 123.0000 & 1820.3034 & 1371.6415 & 0.0000 & 43.0664 & 0.0000 & 0.0000 & n/a & 0.0000 & 2.0000 & 0.6053 \\
CPU bursts & HPA CPU-only (CPU baseline) & 1 & yes & 0.5250 & 130.0000 & 1703.0794 & 2304.0822 & 0.0000 & 43.9414 & 0.0000 & 0.0000 & n/a & 0.0000 & 0.0000 & 0.6250 \\
CPU bursts & PPO (multi-metric, 11 signals) & 1 & yes & 0.5000 & 84.0000 & 1567.6466 & 1263.5172 & 0.0000 & 44.6680 & 3.0000 & 0.0000 & n/a & 0.0000 & 0.0000 & 0.7750 \\
memory ramp & HPA CPU+Memory (resource baseline) & 1 & yes & 0.0000 & 78.0000 & 975.0000 & 975.0000 & 0.0000 & 119.4609 & 0.0000 & 0.0000 & n/a & 0.0000 & 0.0000 & 1.0000 \\
memory ramp & HPA CPU-only (CPU baseline) & 1 & yes & 0.0000 & 40.0000 & 975.0000 & 975.0000 & 0.0000 & 152.3359 & 0.0000 & 0.0000 & n/a & 0.0000 & 0.0000 & 1.0000 \\
memory ramp & PPO (multi-metric, 11 signals) & 1 & yes & 0.0000 & 56.0000 & 950.6250 & 975.0000 & 0.0000 & 120.0742 & 1.0000 & 0.0000 & n/a & 0.0000 & 0.0000 & 0.9750 \\
mixed bursts & HPA CPU+Memory (resource baseline) & 1 & yes & 0.0000 & 76.0000 & 980.3134 & 978.8153 & 0.0000 & 88.1680 & 0.0000 & 0.0000 & n/a & 0.0000 & 1.0000 & 1.0000 \\
mixed bursts & HPA CPU-only (CPU baseline) & 1 & yes & 0.1000 & 62.5000 & 1041.4198 & 982.3931 & 0.0000 & 103.7578 & 0.0000 & 0.0000 & n/a & 0.0000 & 0.0000 & 1.0000 \\
mixed bursts & PPO (multi-metric, 11 signals) & 1 & yes & 0.2000 & 66.5000 & 1091.7229 & 983.7270 & 0.0000 & 87.7266 & 2.0000 & 0.0000 & n/a & 0.0000 & 0.0000 & 0.9750 \\
\hline
\end{tabular}
}
\end{table}
\end{landscape}
